\begin{table*}[h]
\centering
\caption{Mean absolute error (MAE) of benchmark prediction methods under varying query budgets for four HELM-Lite tasks. Results averaged over 1024 random query subsets of size $m$ using all available reference models for DKPS and Ensemble. 
Evaluation is under the Leave-One-Family-Out protocol.
Lower is better; bold indicates lowest MAE per (task, $m$) pair. DKPS-based methods consistently outperform baselines at low query budgets, achieving comparable accuracy with substantially fewer queries (e.g., on legalbench, DKPS at $m=1$ achieves MAE $\approx$ Sample Score at $m=16$, a 16$\times$ reduction).}
\small
\setlength{\tabcolsep}{2.25pt}
\begin{tabular}{l *{16}{S[table-format=1.3]}}
\toprule
& \multicolumn{4}{c}{\textbf{legalbench}} & \multicolumn{4}{c}{\textbf{medqa}} & \multicolumn{4}{c}{\textbf{wmt\_14}} & \multicolumn{4}{c}{\textbf{math}} \\
\cmidrule(lr){2-5} \cmidrule(lr){6-9} \cmidrule(lr){10-13} \cmidrule(lr){14-17}
\textit{Num. queries} & {1} & {4} & {16} & {64} & {1} & {4} & {16} & {64} & {1} & {4} & {16} & {64} & {1} & {4} & {16} & {64} \\
\midrule
Pop. Mean & 0.093 & 0.093 & 0.093 & 0.093 & 0.131 & 0.131 & 0.131 & 0.131 & 0.048 & 0.048 & 0.048 & 0.048 & 0.234 & 0.234 & 0.234 & 0.234 \\
Sample Score & 0.472 & 0.199 & 0.097 & 0.048 & 0.420 & 0.189 & 0.093 & 0.044 & 0.153 & 0.073 & 0.037 & 0.017 & 0.339 & 0.163 & 0.079 & \textbf{0.037} \\
DKPS & \textbf{0.087} & \textbf{0.067} & 0.051 & 0.034 & \textbf{0.121} & \textbf{0.094} & 0.066 & 0.039 & \textbf{0.038} & \textbf{0.026} & 0.019 & 0.016 & 0.144 & 0.105 & 0.082 & 0.070 \\
Ensemble & \textbf{0.087} & \textbf{0.067} & \textbf{0.049} & \textbf{0.033} & \textbf{0.121} & \textbf{0.094} & \textbf{0.065} & \textbf{0.038} & \textbf{0.038} & \textbf{0.026} & \textbf{0.018} & \textbf{0.013} & \textbf{0.143} & \textbf{0.104} & \textbf{0.074} & 0.042 \\
\bottomrule
\end{tabular}
\label{tab:dkps}
\end{table*}
% }

\section{Introduction}
Benchmarks for machine learning methods are widely used to measure the field’s progress.
While they are not without flaw \citep{pmlr-v97-recht19a}, the ubiquity and longevity of early benchmarks such as MNIST \citep{lecun2010mnist}, CIFAR-10 \citep{Krizhevsky09learningmultiple}, and the Iris dataset \citep{fisher36lda} attest to their practical utility.
In the era of modern generative models, benchmarks have shifted away from evaluating performance on a single, well-defined task and moved towards evaluating performance across a suite of tasks designed to capture a model’s overall capability. 
For example, \citep{liang2022holistic} introduced the HELM benchmark suite and leaderboard as an attempt to evaluate language models as comprehensively and transparently as possible. 
Similar multi-task benchmark suites and leaderboards now exist for embedding models \citep{muennighoff2023mtebmassivetextembedding}, image generation \citep{huang2025t2icompbenchenhancedcomprehensivebenchmark}, coding \citep{chen2021evaluatinglargelanguagemodels}, and ``intelligence" \citep{chollet2019measureintelligence}, among others.

Although multi-task evaluations offer a more complete picture of a model’s strengths and weaknesses, they are often computationally expensive. 
Evaluating a model on a modern multi-task benchmark typically requires generating and scoring thousands of responses, making full evaluation increasingly impractical as both model sizes and benchmark scopes grow. 
At the same time, it is easier than ever to produce new model behavior -- via parameter efficient fine-tuning \citep{han2024parameterefficientfinetuninglargemodels}, instruction tuning \citep{zhang2025instructiontuninglargelanguage}, model merging \citep{model-merging}, distillation \citep{Gou_2021}, prompt engineering \citep{sahoo2025systematicsurveypromptengineering}, etc. -- making it infeasible to evaluate each new variant.

These trends motivate the need for query-efficient alternatives to full benchmark evaluation. 
In this paper, we address this problem by leveraging information and responses from previously evaluated models under the assumption that it is possible to predict a new model’s score by leveraging the similarity between its responses and cached responses from already-scored models on a (small) subset of queries. 

\noindent \textbf{Contribution.} We make two primary contributions.
\begin{enumerate}[itemsep=0pt, topsep=0pt, parsep=2pt]
\item \textbf{Theoretical:} We prove that a benchmark prediction method based on the Data Kernel Perspective Space (DKPS) is query-efficient relative to estimating benchmark performance using a subset of queries.

\item \textbf{Empirical:} We validate theoretical query-efficiency of DKPS-based methods. 
As summarized in Table~\ref{tab:dkps}, DKPS-based methods can reduce the number of queries required to achieve a given performance by more than 10$\times$.
\end{enumerate}
\subsection{Background \& related work}


\paragraph{Low-dimensional representations of language models.}
Our work contributes to a growing literature on low-dimensional representations of language models. 
Some approaches construct representations by comparing internal activations across models \citep{duderstadt2023comparing, huh2024platonicrepresentationhypothesis, horwitz2025chartatlasworldsmodels} or by comparing model weights directly \citep{chen2025extracting}. 
However, evaluation frameworks typically do not require (or even allow) access to model internals, and instead rely only on model responses.
The Data Kernel Perspective Space (DKPS), first introduced in the context of monitoring multi-agent systems \citep{helm2024tracking}, addresses this setting by mapping a collection of generative models into a low-dimensional Euclidean space via multidimensional scaling of matrices containing average embedded responses. 
Recent theoretical work has shown that DKPS recovers ground-truth low-dimensional model structure as the number of models, queries, and responses per query grow \citep{aranyak,acharyya2025concentrationboundsresponsebasedvector} and that supervised inference methods trained on DKPS representations are consistent \citep{helm-etal-2025-statistical}.
We frame predicting the score of a model on a benchmark as an instance of model-level inference in the black-box setting and extend recent theoretical and empirical results realted to the DKPS to query-efficiency for benchmark score prediction.

\paragraph{Efficient benchmarking.}

Several lines of work aim to reduce the computational cost of benchmark evaluation.
\citet{polo2024tinybenchmarks} use Item Response Theory to construct benchmark subsets of 100 examples (approx. 1\% of original size) that predict performance within 2\% error.
\citet{vivek-etal-2024-anchor} propose clustering examples based on cross-model predictions and selecting cluster centers as "anchor points," achieving accurate rankings with significantly fewer examples.
\citet{bean2025scales} propose item-centric selection based on cognitive embeddings, improving cold-start performance and cross-family transferability.
\citet{li2024active} use reinforcement learning to model dependencies across examples, reducing estimation error by 25-50\% via active selection.
\citet{perlitz2024efficient} analyze benchmark design choices and propose ranking algorithms achieving 100× cost reductions with minimal ``reliability" loss.
These techniques typically require explicit or implicit task-specific structure, model metadata, access to model internals, or an response-level scoring functions.

Our approach is complementary to prior work in that we leverage cached responses from previously evaluated models, operate purely via black-box access and model response similarity, and assume no structure across tasks or subtasks. 
Importantly, this enables combining our method with existing techniques -- e.g., using DKPS-based methods to predict performance on IRT-selected subsets \citep{polo2024tinybenchmarks} or on  ``anchor points" \citep{vivek-etal-2024-anchor} -- to potentially compound efficiency gains beyond what either approach achieves alone.

\subsection{Problem statement}
Given a generative model $ f $, we consider the problem of predicting its score on a benchmark $ Q^{*} = \{q_{1}, \hdots, q_{M}\} $.
Let $ y: \mathcal{F} \times 2^{Q^{*}} \to [0,1] $ denote the benchmark scoring function.
That is, $ y(f, Q^{*}) $ is the score assigned to $ f $ after evaluating it on all queries.
We refer to $ y(f, Q^{*})$ as $ y $ when the context is clear.

As described above, full evaluation of $ f $ on all queries may be infeasible. 
Instead, we assume access to previously-evaluated (model, score) pairs $ (f_{1}, y_{1}), \hdots (f_{n}, y_{n}) $ and each model's response to each query $ \{\{f_{i}(q_{j})\}_{j=1}^{M}\}_{i=1}^{n} $.
Given this additional information, our goal is to estimate $ y $ with $ m \ll M $ queries.
That is, we seek to identify an estimate $ \hat{y} $ that predicts the full benchmark score of $
f $ from its responses to a subset $ Q \in 2{^{Q^*}} $ of the benchmark queries and the information available from the previously-evaluated models.